\clearpage
\unitheader{C}{How does batch size affect the hyperparameters?}{}

We now have had some taste of bivariate hyperparameters scaling (in our learning rate and weight decay tuning problem). We might wonder whether if there is a better way, where we can mathematically derive the right scaling relationships between hyperparameters.

To this end, we will study how batch size, learning rate, and scale all interact with each other. Batch size is an ideal variable to study, since the relationship between batch size, learning rates and convergence rates are very well studied and understood as part of the theory of online optimization.

We will make use of a well-known toy model called the noisy quadratic model, and see how this gives us learning-rate rules that (might) carry over to language model training in the wild.

\begingroup
\setcounter{question}{0}
\renewcommand{\aTwoProblemNumber}{3.\arabic{question}}

\subunitheader{Background: The noisy quadratic model}

A noisy quadratic model (NQM) replaces neural-network training with a quadratic
loss and noisy gradients. There are two key properties of the noisy quadratic that we will need to vary and understand. The curvature of the quadratic controls the optimization difficulty of the problem, and controls how quickly different directions
can be optimized; the gradient noise represents variation among training examples and informs whether larger batches are helpful.
We will follow
\href{https://proceedings.neurips.cc/paper_files/paper/2019/hash/e0eacd983971634327ae1819ea8b6214-Abstract.html}{Guodong Zhang et al. (2019),
\emph{Which Algorithmic Choices Matter at Which Batch Sizes? Insights From a
Noisy Quadratic Model}}, which uses this simple model to explain how learning
rate, momentum, and the optimizer affect batch-size scaling.

We will investigate two conclusions and the conditions under which they hold:
\begin{enumerate}[label=\arabic*.]
\item \textbf{Optimal LR can follow a power law in batch size.}
Let $B$ be batch size and $\eta^*(B)$ the LR minimizing expected final loss
at a fixed total number of training examples. Write a candidate rule as
$\eta^*(B)\approx cB^p$, with fitted constants $c$ and $p$.
For SGD in the small-step regime, the NQM predicts approximately linear scaling,
        $p=1$: doubling batch size calls for doubling LR. Of course, there will come a point where this prediction breaks down, as an arbitrarily large learning rate will lead to divergence. We will test this scaling rule, and identify the range of parameters where it is useful.
\item \textbf{Momentum can help more at large batches than at small batches.}
The classical NQM prediction is that, in the noise-dominated small-batch regime,
SGD momentum mainly rescales the effective LR and offers little extra benefit
after LR tuning. At larger batches, reduced noise can allow momentum to
accelerate convergence. We will test this prediction at a finite training budget and perform comparisons with Adam.
\end{enumerate}


Our first problem will make use of a simple, two dimensional quadratic model for you to play with.
\paragraph{2D quadratic toy model.}
Let $w=(w_1,w_2)^\top$ be the parameters and fix the curvature matrix
$H=\operatorname{diag}(1,10)$. The loss and its exact gradient are
\[
 f(w)=\tfrac12 w^\top H w=\tfrac12(w_1^2+10w_2^2),
 \qquad \nabla f(w)=Hw.
\]
To define the \emph{noisy} part of this quadratic, we will consider our gradients to be noisy version of the true gradient. For an update with batch size $B$, we define the corresponding minibatch gradient to be
\[
 g_t=Hw_t+\epsilon_t,\qquad
 \epsilon_t\sim\mathcal N(0,I_2/B)
\]
independently across updates, where $I_2$ is the two-dimensional identity matrix. This noise model comes from the naive assumption that the gradients are averages of i.i.d. items, and thus we expect variance to scale as $O(1/B)$.

\paragraph{Optimizers.}
In this problem, we will use a simplified optimizer model with constant LR $\eta$ and no weight decay. For the optimizer itself, we will study Stochastic gradient descent (SGD)  (defined as $w_{t+1}=w_t-\eta g_t$), Adam with $\epsilon=10^{-8}$, $\beta_{1}=0.9$, and $\beta_{2}=0.95$, and RMSProp, which is Adam with $\beta_{1}=0$.

\begin{problem}{What does this model predict about batch size? (Run locally with CPUs)}
  Consider our 2D toy quadratic model above. You will perform stochastic optimization on this toy model using the noisy gradients and compare how the hyperparameters change the final loss.

    Concretely, let $N=8,192$ be the total dataset size, $B$ be the batch size (which we specify below), and  $w_{0,1}\sim\mathcal N(0,1)$ and $w_{0,2}\sim\mathcal N(0,.1)$ as initialization.

     Compare expected final loss $L_N(B,\eta)=\mathbb E[f(w_{N/B})]$, estimated by averaging over independent
     initializations and noise (using at least 512 samples).

  \begin{parts}


\item \textbf{SGD.}
Sweep LR at batches $B\in\{1,2,4,8,16,32,64\}$ and fit a power law to the
estimated optimal LRs. Predict the optimal LR at $B=256$, then test the
prediction with an LR sweep at that batch. Extend the batch range through
512 and plot optimal LR and best final loss after tuning LR.
Over what range does the rule work? Think about how reduced gradient noise
competes with fewer updates at fixed $N$.

\item \textbf{Adaptive updates.}
Repeat the experiment with RMSProp and Adam, tuning LR separately.
Compare the source exponents and the losses at predicted versus tuned LRs
at batch 256. Do these optimizers follow the same rule, and does a similar exponent imply similar transfer accuracy?
State your predictions for the language model setting. How should optimal LR change with batch size, and over what range do you expect that prediction to hold?

\item \textbf{Explore the assumptions.}
Replace the gradient-noise covariance $I_2$ by $\sigma^2 I_2$ for
$\sigma\in\{1,\allowbreak 10,\allowbreak 100,\allowbreak 300\}$. Repeat the RMSProp and Adam source fits and
batch-256 tests. Plot the fitted exponent against $\sigma$.
Does changing the ratio of gradient signal to noise change the scaling rule?
Relate your observations to the contributions to the second moment.

As a curvature comparison, repeat the experiments with the scalar loss $f(w)=w^2/2$, initial $w_0\sim\mathcal N(0,2)$, and gradient-noise variance $\sigma^2/B$. Which conclusions
are sensitive to the curvature model?

\item \textbf{Momentum at small and large batches.}
Return to $H=\operatorname{diag}(1,10)$ and noise covariance $I_2$.
For SGD with momentum, initialize $s_0=0$ and use
\[
 s_{t+1}=\mu s_t+g_t,\qquad w_{t+1}=w_t-\eta s_{t+1},
\]
where $\mu$ is the momentum coefficient.
At small batch $B=16$ and large batch $B=256$, compare $\mu=0$ and $.9$,
tuning LR separately. For Adam, keep $\beta_2=.95$ and jointly sweep
$\beta_1$ and LR at each batch (hint: you should include $\beta_{1}=0$, $0.9$ and other values closer to one).

Plot the best loss at each $\beta_1$ after tuning LR.
How does the benefit of momentum depend on batch size? Compare the learning
curves of the best ($\beta$, LR) pair with the best no-momentum baseline. Do you expect these observations to carry over to language model training?

\end{parts}
\end{problem}

\clearpage
\subunitheader{Batch size scaling in a language model}

Do the scaling rules and momentum effects from the NQM carry over to
language-model training? Use the default model and training settings at
614.4M tokens, varying batch size and the specified hyperparameters.

Here $B$ is the number of sequences per optimizer update. Use microbatches of
$\min(B,64)$ sequences and report the actual number of tokens processed.
Use a linear learning-rate schedule with warmup over the first 1\% of
optimizer updates. The batch-switch example below specifies its own
schedule in terms of tokens processed.


\begin{problem}{What should scale when batch size changes? (suggested: 48 new runs)}
\begin{parts}
\item \textbf{Tune LR with weight decay (9 new runs).}
Set WD to $.1$ and sweep peak LR at $B\in\{8,16,32,64\}$.
Plot the loss--LR curves, fitted optimal LR, and best measured loss against
batch size. How do the optimal LR and performance change with batch size?

\item \textbf{Compare two scaling hypotheses (suggested: 27 new runs).}
The NQM can also give predictions of how weight decay and batch might jointly scale together.
Design experiments at $B\in\{8,16,32,64\}$ to test two hypotheses:
\begin{enumerate}[label=\roman*.]
\item Fix WD at $.1$ across batches and scale the optimal peak LR.
\item Fix peak LR at $.0015$ across batches and scale the optimal WD.
\end{enumerate}
Use the source results to predict configurations at $B=128$ and $B=256$,
then test which hypothesis transfers better. Discuss whether your findings
agree with the NQM.

\item \textbf{Ablate momentum (suggested: 12 new runs).}
At $B=8$ and $B=256$, hold each batch's best measured LR--WD pair and
$\beta_2=.95$ fixed. Design experiments including a $\beta_1=0$ control to test how the effect of momentum changes with batch size.
Compare with our earlier problem when we retuned LR.

\item \textbf{What does the NQM explain? (no new runs).}
Which predictions from your quadratic experiments agree with your language-model experiments,
and which do not?

Give a short prediction--measurement table for LR--batch-size scaling,
weight decay, and momentum. Include the measured target losses supporting
your conclusions. Distinguish a prediction that fails from an effect
the NQM does not model. What would you change in the NQM to explain one of
the disagreements?
\end{parts}
\end{problem}

\paragraph{Further reading.}
\href{https://arxiv.org/abs/2505.13738}{\textbf{Power Lines}}: WD--LR--batch coupling;
\href{https://arxiv.org/abs/2503.04715}{\textbf{Step Law}}: optimal LR and batch-size scaling;
\href{https://arxiv.org/abs/2606.16899}{\textbf{Hyperball}}: control of weight and update norms.

\endgroup
\setcounter{question}{3}

\clearpage
\begin{examplequestion}[difficulty=hard]
\phantomsection\label{ex:a2-batch-switch}
\textbf{Double batchsize during training}
Train the $d8$ language model 614.4M tokens with AdamW,
initial batch size $64$, peak LR $.0015$, and WD $1.6$.
Let $\eta_{\rm base}(D)$ denote the learning rate after processing
$D$ tokens: warm up over the first $1\%$ of tokens, then decay linearly to zero
at 614.4M tokens.

After processing exactly 307.2M tokens, double the batch size to $128$.
Reach the midpoint with 4,687 batches of 64 sequences followed by one batch
of 32; this gives 4,688 updates before the switch. Then use 2,343 batches of
128 sequences and a final batch of 96, giving 2,344 updates after the switch.
Compute each update's mean loss using its actual batch size.
Keep Adam's accumulated moments and both $\beta$ values unchanged.
For the remaining half of training, consider five policies:
\begin{center}
\begin{tabular}{clcc}
\toprule
Policy & Change at the midpoint & LR after the change & WD after the change\\
\midrule
A & Neither LR nor WD & $\eta_{\rm base}(D)$ & $1.6$\\
B & Multiply LR by $\sqrt{2}$ & $\sqrt{2}\,\eta_{\rm base}(D)$ & $1.6$\\
C & Multiply LR by $2$ & $2\,\eta_{\rm base}(D)$ & $1.6$\\
D & Multiply WD by $\sqrt{2}$ & $\eta_{\rm base}(D)$ & $1.6\sqrt{2}$\\
E & Multiply WD by $2$ & $\eta_{\rm base}(D)$ & $3.2$\\
\bottomrule
\end{tabular}
\end{center}
All runs process the same ordered tokens. Evaluate $\eta_{\rm base}(D)$ at
the number of tokens consumed \emph{before} each update, including the partial
batches. Changing the batch size must not restart warmup.
\textbf{Which policy do you predict will give the lowest mean final validation loss across 3 seeds (42,43,44)?}
Choose one and explain how changing LR or WD affects training when there are approximately half as many optimizer updates per token after the switch.
\end{examplequestion}

